% Copy this file to auction/MM-DD-YY/main.tex.

\documentclass[../../main.tex]{subfiles}

\begin{document}

\section{Motivations}

Our goal is to use geometry to visualize bidder preferences. Here is an example:

\begin{example} \label{graph:sub-comp}
    A hotel offers exactly two rooms: room 1 and room 2. Denote room 1's price as $p_1$ and room 2's price as $p_2$. Sunny wants no more than one room. His utility function can be described as
    \[
        U_{\text{Sunny}}(0,0)=0,\quad
        U_{\text{Sunny}}(1,0)=30-p_1,\quad
        U_{\text{Sunny}}(0,1)=40-p_2.
    \]

    We can visualize Sunny's preferences in the following graph.

    % Figure for Sunny

    Note that room 1 and room 2 are substitutes for Sunny.

    Rainy, on the other hand, wants both rooms or no room. Her utility function can be described as
    \[
        U_{\text{Rainy}}(0,0)=0,\quad
        U_{\text{Rainy}}(1,1)=50-p_1-p_2.
    \]

    Similarly, we can visualize Rainy's preferences in the following graph.

    % Figure for Rainy

    Note that room 1 and room 2 are complements for Rainy.
\end{example}

We say a \textit{competitive equilibrium} exists at price $p$ if there is a complete assignment of the supply to all agents such that each agent receives a bundle that maximizes its utility.

\begin{example}
    In the previous example \ref{graph:sub-comp}, we can visualize both Sunny and Rainy's demands on one graph by superposing figures %\ref{} here.
    % Figure for both
    
    In this case, since the supply is exactly one copy of each room, i.e., $s=(1,1)$, a competitive equilibrium does not exist.

    However, suppose Rainy really loves the hotel and is willing to pay more. To be clear, her updated utility function becomes
    \[
        U'_{\text{Rainy}}(0,0)=0,\quad
        U'_{\text{Rainy}}(1,1)=80-p_1-p_2.
    \]

    % Updated figure
    If we look at the updated graph, a competitive equilibrium exists in the shaded triangle.
\end{example}

The rest of these notes develops the geometry underlying these pictures. We first study the geometry generated by a single agent, then understand how these geometries combine across agents, and finally use this structure to study existence and counting of competitive equilibria.

\section{Demand Geometry of a Single Agent}

\subsection{Locus of indifference prices}

Let $J$ be the set of all agents. For each agent $j \in J$, let $$X^j \subseteq \zz^n$$ be bundles available to $j$. The agent $j$'s valuation is a function $$V^j:X^j \to \rr,$$ and it has a quasi-linear utility function $V^j(x)-p\cdot x$, where $p\in \rr^n$ and $x \in X^j$. Hence, agent $j$ demands bundles in 
\[D_j(p)=\argmax_{x \in X^j} \{V^j(x)-p\cdot x\}\]

For a fixed bundle $x$, 
\[V^j(x)-p\cdot x\]
is affine-linear in $p$. Therefore, the price space is divided into regions according to which affine function achieves the maximum.

Given these notations, the competitive equilibrium can be defined in the following way.

\begin{defin}
    A \textit{competitive equilibrium} at a price $p$ is a state of a market or economy with supply $s$ where 
    \[x^j \in D^j(p)\text{ for every }j \in J,\]
    and
    \[\sum_{j \in J}x^j=s.\]
\end{defin}

\begin{defin}
    The \textit{Locus of Indifference Price (LIP)} is the set of prices at which the agent is indifferent between more than one demanded bundels. Mathematically, 
    \[L_j := \{p \in \rr^n:|D_j(p)|>1\}.\]

    Note that the linear pieces in previous figures %ref
    are graphs of the LIP.
\end{defin}

% facet, lower-dimensional intersections

\subsection{Demand complex}
\subsection{Facets and bundle differences}
\subsection{Demand types}

\section{Aggregating Agents}

\subsection{Aggregate demand}
\subsection{Minkowski sums}
\subsection{Competitive equilibrium as market clearing}
\subsection{Holes and discrete convexity}

\section{Unimodularity}

\subsection{Lattice intuition}
\subsection{Unimodular demand types}
\subsection{The Unimodularity Theorem}

\section{Intersections and Equilibria}

\subsection{Superposing demand complexes}
\subsection{Intersection multiplicity}
\subsection{Intersection Counting Theorem}

\end{document}
